\documentclass[10pt,a4paper]{article}

\usepackage[utf8]{inputenc}
\usepackage[T1]{fontenc}
\usepackage{lmodern}
\usepackage[margin=1in]{geometry}
\usepackage{amsmath,amssymb,amsthm,mathtools}
\usepackage{booktabs}
\usepackage{array}
\usepackage{enumitem}
\usepackage{xcolor}
\usepackage{tikz}
\usepackage{placeins}
\usepackage{hyperref}

\hypersetup{
  colorlinks=true,
  linkcolor=blue!60!black,
  urlcolor=blue!60!black,
  citecolor=blue!60!black
}

\newtheorem{theorem}{Theorem}[section]
\newtheorem{lemma}[theorem]{Lemma}
\newtheorem{proposition}[theorem]{Proposition}
\newtheorem{definition}[theorem]{Definition}
\theoremstyle{remark}
\newtheorem{remark}[theorem]{Remark}
\newtheorem{example}[theorem]{Example}

\newcommand{\val}{\operatorname{val}}

\setlength{\parindent}{0pt}
\setlength{\parskip}{6pt plus 2pt minus 1pt}

% \newcommand{\cl}[1]{[#1]}
% \newcommand{\card}[1]{\lvert #1 \rvert}

\title{\textbf{
  Cyclic $n \!\times\! n$ Latin squares \\ 
  representing solutions of \\
  $(n,c)$-coins problems for $0 \le c \le n$
}}
\author{Marc Duponcheel and Luc Duponcheel}

\date{\today}

\begin{document}

\maketitle

\begin{abstract}

\noindent The \emph{$n$-queens problem} on an $n \!\times\! n$ square allows at most $1$ queen on any
row, column, sum-diagonal, or diff-diagonal ($n = 1$ is trivial and $n \in \{2,3\}$ have no solution,
so \emph{we limit ourselves to $n \ge 4$}). The \emph{$(n,c)$-coins problem} ($0 \le c \le n$)
generalizes $n$-queens by allowing at most $c$ coins on any row, column, sum-diagonal, or
diff-diagonal ($c = 1$ is $n$-queens). \emph{We prove}:
\begin{itemize}[nosep]
\item For every $n \ge 4$, an \emph{incremental $n \!\times\! n$ Latin square} represents $n + 1$
nested solutions $S^n_0 \subset S^n_1 \subset \dots \subset S^n_n$ of $(n,c)$-coins problems
\emph{simultaneously} for all $0 \le c \le n$ (Theorem~\ref{thm:main1}).
\item For every $n \ge 4$ with $n \notin \{4, 6\}$, there \emph{exists} an incremental \emph{cyclic}
$n \!\times\! n$ Latin square, given by four infinite modular families modulo~$6$ for
$n \notin \{4, 6, 8, 9, 15, 21\}$ and explicit cyclic constructions for $n \in \{8, 9, 15, 21\}$
(Theorem~\ref{thm:main2}, Appendix~\ref{app:small}, and~\cite{PalomarLean,CompanionRepo}).
\item For $n \in \{4, 6\}$, no \emph{cyclic} incremental Latin square exists (by exhaustive search
over all $4! = 24$ and $6! = 720$ permutations), whereas explicit \emph{non-cyclic} incremental Latin
squares $L^4, L^6$ exist (Appendix~\ref{app:noncyclic} and~\cite{PalomarLean,CompanionRepo}); thus an
incremental $n \!\times\! n$ Latin square exists for every $n \ge 4$ (Theorems~\ref{thm:main3}--\ref{thm:main4}).
\end{itemize}

\noindent \emph{Key methodology}:
\begin{itemize}[nosep]
\item Whereas Le, Li and Wang~\cite{LLW1990} construct $(n,c)$-coins solutions \emph{separately} for
each fixed $c$, an incremental Latin square partitions the board into $n$ disjoint transversals whose
sublevel sets solve all $0 \le c \le n$ \emph{simultaneously} (Section~2).
\item The \emph{injective-majorization principle} (Lemma~\ref{lem:inj-maj}, Section~3) replaces
$O(n^3)$ naive diagonal prefix verification by a single majorized injection per diagonal direction.
\item All $n \notin \{4, 6, 8, 9, 15, 21\}$ are partitioned into four infinite families modulo~$6$
(Section~4): (1)~$n = 1,5 \bmod 6$ ($n \ge 5$), (2)~$n = 2,4 \bmod 6$ ($n \ge 10$),
(3)~$n = 3 \bmod 6$ ($n \ge 27$), and (4)~$n = 0 \bmod 6$ ($n \ge 12$), each $O(1)$-time per entry
($O(n^2)$ for the full square vs.\ $\Theta(n^3)$ to generate all separate solutions in~\cite{LLW1990}). All
results (including $n \in \{8, 9, 15, 21\}$ and $n \in \{4, 6\}$) are mechanized in
Lean~4~\cite{PalomarLean,CompanionRepo}.
\end{itemize}

\noindent Repositories and discussions:
\begin{itemize}[nosep]
\item Companion repo~\cite{CompanionRepo}: {\scriptsize\url{https://github.com/NofanTasi/incremental-nc-coins-problem}}
\item Palomar Lean~4~\cite{PalomarLean}: {\scriptsize\url{https://github.com/LucDuponcheelAtGitHub/incremental-nc-coins}}
\item Discussions: \href{https://mathoverflow.net/questions/510135}{MO \#510135}~\cite{MO510135}, \href{https://puzzling.stackexchange.com/questions/138939}{PSE \#138939}~\cite{PSE138939}, \href{https://mathoverflow.net/questions/516225}{MO \#516225}~\cite{MO516225}.
\end{itemize}

\end{abstract}

\section{Assumptions, Notation and Definitions}

Throughout, $n \in \mathbb{N}$ satisfies $n \ge 4$, $[m,n] := \{m, \dots, n\}$ ($m \le n$), and
$S^n : [0,n-1]^2$ is a $0$-based $n \!\times\! n$ square. We use \emph{uppercase} for cell sets and
\emph{lowercase} for their values:
\begin{itemize}[nosep]
\item $S^n_{i,j}$ ($s^n_{i,j}$): (the value of) the $(i,j)$-cell of $S^n$ ($i,j \in [0,n-1]$).
\item $R^n_i$ ($r^n_{i,j}$): $\{S^n_{i,j} \mid j \in [0,n-1]\}$, the $i$-row of $S^n$ ($i \in [0,n-1]$).
\item $C^n_j$ ($c^n_{i,j}$): $\{S^n_{i,j} \mid i \in [0,n-1]\}$, the $j$-column of $S^n$ ($j \in [0,n-1]$).
\item $D^n_{i+j=s}$ ($d^n_{i+j=s}$): $\{S^n_{i,j} \mid i + j = s\}$, the $s$-sum-diagonal ($s \in [0,2(n-1)]$).
\item $D^n_{i-j=d}$ ($d^n_{i-j=d}$): $\{S^n_{i,j} \mid i - j = d\}$, the $d$-diff-diagonal ($d \in [-(n-1),n-1]$).
\end{itemize}

\begin{definition}\label{def:nc}
A solution of an $(n,c)$-coins problem ($c \in [0,n]$) is a subset $S^n_c \subseteq S^n$ of $n\,.\,c$
cells such that every row $R^n_i$, column $C^n_j$, sum-diagonal $D^n_{i+j=s}$, and diff-diagonal
$D^n_{i-j=d}$ contains at most $c$ cells of $S^n_c$.
\end{definition}

\begin{definition}\label{def:is}
For each threshold $k \in [0,n-1]$, define the diagonal sublevel sets
\[
    D^{n,k}_{i+j=s} := \{S^n_{i,j} \in D^n_{i+j=s} \mid s^n_{i,j} \le k\}
    \;\;\text{and}\;\;
    D^{n,k}_{i-j=d} := \{S^n_{i,j} \in D^n_{i-j=d} \mid s^n_{i,j} \le k\}.
\]
A square $S^n$ is \emph{incremental} iff for every $k \in [0,n-1]$,
$|\,D^{n,k}_{i+j=s}\,| \le k + 1$ for all $s \in [0,2(n-1)]$ and $|\,D^{n,k}_{i-j=d}\,| \le k + 1$
for all $d \in [-(n-1),n-1]$.
\end{definition}

\begin{remark}\label{rem:row-col}
Definition~\ref{def:is} explicitly restricts only diagonal sublevel sets. In any Latin square $L^n$, each row and column is a permutation of $[0,n-1]$, so the row and column sublevel sets $\{L^n_{i,j} \in R^n_i \mid l^n_{i,j} \le k\}$ and $\{L^n_{i,j} \in C^n_j \mid l^n_{i,j} \le k\}$ automatically contain exactly $k + 1$ cells for every $k \in [0,n-1]$. Thus row and column capacity constraints are inherently satisfied whenever $S^n$ is a Latin square.
\end{remark}

\begin{definition}\label{def:ls}
A square $S^n$ is a \emph{Latin} square ($L^n$, with values $l^n_{i,j}$) iff every row $r^n_i$ and
column $c^n_j$ is a permutation of $[0,n-1]$.
\end{definition}

\begin{definition}\label{def:cs}
A square $S^n$ is \emph{cyclic} iff $s^n_{i + 1, j} = (s^n_{i,j} + 1) \bmod n$ for all
$i,j \in [0,n-1]$. Given a permutation $p$ of $[0,n-1]$, $L^{n,p}$ denotes the cyclic Latin square
with values $l^{n,p}_{i,j} = (i + p(j)) \bmod n$.
\end{definition}

\subsection{Related Work and Novelty}\label{sec:related}

\emph{Survey of related work}: Bell and Stevens~\cite{BS2009} survey $n$-queens results, including
modular/toroidal queens and Latin square transversals. By the Hall--Paige condition, the unpermuted
Cayley table of $\mathbb{Z}_n$ decomposes into transversals iff $n$ is odd, and into toroidal $n$-queens
transversals iff $\gcd(n,6) = 1$~\cite{BS2009,CW2010}.

\emph{Comparison with Le, Li and Wang~\cite{LLW1990}}: Le, Li and Wang introduced the $(n,c)$-coins
problem, constructing solutions $S^n_c$ \emph{separately} for each fixed $c \in [0,n]$ through an
intricate case analysis. In contrast, an \emph{incremental} Latin square $L^n$ partitions the board
into $n$ disjoint transversals whose sublevel sets solve all $c \in [0,n]$ \emph{simultaneously}
(Theorem~\ref{thm:main1}). Generating all $n + 1$ solutions via~\cite{LLW1990} requires
$\sum_{c=1}^n O(n \cdot c) = \Theta(n^3)$ coin placements, whereas our single square provides all
solutions in $O(n^2)$ time ($O(1)$ time per entry, and $O(n)$ space for the permutation $p$).

\emph{Transversal decompositions vs.\ incremental squares}: In classical design theory, a Latin square
partitions the grid into $n$ disjoint transversals (each having $1$ cell per row and column). However,
classical transversals impose \emph{no constraints} on diagonal prefixes. An \emph{incremental} Latin
square requires that every prefix union $\bigcup_{k < c} T_k$ contains at most $c$ cells on \emph{every}
diagonal. This nested property provides layered conflict avoidance (e.g., in multi-priority resource
scheduling, VLSI channel routing, and collision-free memory interleavers for parallel turbo decoders)
where capacity can be expanded dynamically from $c$ to $c+1$ without moving existing coins.

\begin{table}[htbp]
\centering
\scriptsize
\begin{tabular}{@{}llllll@{}}
\toprule
Approach & Board Structure & Range of $n$ & Simultaneous? & Time (all $c$) & Formal Proof \\
\midrule
Le, Li, Wang~\cite{LLW1990} & Separate $S^n_c$ sets & All $n \ge 4$ & No (per-$c$) & $\Theta(n^3)$ & No \\
Toroidal queens~\cite{BS2009,CW2010} & Cayley table & $\gcd(n,6) = 1$ & No diags bound & N/A & Partial \\
This work & Incremental Latin square & All $n \ge 4$ & Yes ($S^n_c = \{l < c\}$) & $O(n^2)$ & Lean~4 (complete) \\
\bottomrule
\end{tabular}
\caption{Comparison of approaches to $(n,c)$-coins problems and diagonal combinatorial designs.}
\label{tab:comparison}
\end{table}

\emph{Other works on $n$-queens and cyclic Latin squares}:
\begin{itemize}[nosep]
\item Pratt~\cite{Pratt2016} gives permanent formulas counting $n$-queens ($c = 1$) and Latin squares;
it is purely \emph{enumerative} ($c = 1$) and does not construct incremental squares or $(n,c)$-coins
solutions for $c > 1$.
\item Cavenagh and Wanless~\cite{CW2010} bound transversal counts in the \emph{unpermuted} Cayley table
$(i + j) \bmod n$ for odd $n$. By contrast, our cyclic squares permute columns via
$l^{n,p}_{i,j} = (i + p(j)) \bmod n$, exist for both odd and even $n \notin \{4,6\}$, and satisfy
\emph{non-modular} diagonal prefix bounds $|\,D^{n,k}\,| \le k + 1$ for all $k \in [0,n-1]$.
\item Mirzavaziri and Yaqubi~\cite{MY2026} study cyclic Latin Eulerian numbers (column ascents),
unrelated to diagonal coin capacities.
\end{itemize}

\section[Incremental Latin Squares and (n,c)-Coins Problems]{Incremental Latin Squares and $(n,c)$-Coins Problems}

\begin{theorem}\label{thm:main1}
For every $n \ge 4$, an incremental $n \!\times\! n$ Latin square $L^n$ represents $n + 1$ solutions
of $(n,c)$-coins problems \emph{simultaneously} for all $c \in [0,n]$ as
$S^n_c = \{L^n_{i,j} \mid l^n_{i,j} < c\}$ ($S^n_0 = \emptyset$).
\end{theorem}

\begin{proof}
For $c = 0$, $|\,S^n_0\,| = 0 = n\,.\,0$. For $c \in [1,n]$,
$S^n_c = \{L^n_{i,j} \mid l^n_{i,j} \le c - 1\}$. Since each row $r^n_i$ and column $c^n_j$ is a
permutation of $[0,n-1]$, $|\,S^n_c \cap R^n_i\,| = |\,S^n_c \cap C^n_j\,| = |\,[0,c-1]\,| = c$ and
$|\,S^n_c\,| = n\,.\,c$. By Definition~\ref{def:is} with $k = c - 1$,
$|\,S^n_c \cap D^n_{i+j=s}\,| = |\,D^{n,c-1}_{i+j=s}\,| \le c$ and
$|\,S^n_c \cap D^n_{i-j=d}\,| = |\,D^{n,c-1}_{i-j=d}\,| \le c$.
\end{proof}

\section{Injective Majorization}\label{sec:inj-maj}

To avoid $O(n^3)$ naive cardinality checks across all $n$ thresholds and $4n - 2$ diagonals (where
counting cells of value $\le k$ along each diagonal across all $k \in [0,n-1]$ requires
$(4n - 2) \cdot n \cdot O(n) = O(n^3)$ verification operations), the \emph{injective-majorization
principle} reduces Definition~\ref{def:is} on any diagonal $D$ ($D^n_{i+j=s}$ or $D^n_{i-j=d}$) to a
single majorized injection.

\begin{lemma}\label{lem:inj-maj}
If $f : [0,n-1]^2 \to \mathbb{N}$ satisfies (1)~$f(i,j) \le l^{n,p}_{i,j}$ on $D$ and (2)~$f$ is
injective on $D$, then $|\,D^k\,| \le k + 1$ for every threshold $k \in [0,n-1]$.
\end{lemma}

\begin{proof}
By (1), $f$ maps $D^k = \{(i,j) \in D \mid l^{n,p}_{i,j} \le k\}$ into $[0,k]$; by (2), $f|_{D^k}$ is
injective, so $|\,D^k\,| \le |\,[0,k]\,| = k + 1$.
\end{proof}

\section{Incremental Cyclic Latin Squares}

We construct an incremental cyclic Latin square $L^{n,p}$ for every $n \ge 4$ ($n \notin \{4,6\}$) via
four infinite families $F_1$--$F_4$ and four small cases $n \in \{8, 9, 15, 21\}$
(Appendix~\ref{app:small} and~\cite{PalomarLean,CompanionRepo}). Set $m := n/2$ when $2 \mid n$ and
$k := n/3$ when $3 \mid n$. For $x \in \mathbb{Z}$, $x \bmod n \in [0,n-1]$ is the remainder
modulo~$n$. Each modular equality or equivalence is checked in two steps: (i)~\emph{modular reduction}
$(x + qn) \bmod n = x \bmod n$ ($q \in \mathbb{Z}$) to $r \bmod n$, and (ii)~\emph{range verification}
$0 \le r < n$ so $r \bmod n = r$. For brevity, we spell out the intermediate equality and equivalence
steps for Family~$F_2$ as a model and state the resulting reduced forms for $F_3$ and $F_4$ (all steps
are mechanized in~\cite{PalomarLean,CompanionRepo}).

Since $i = s - j$ on $D^n_{i+j=s}$ and $i = d + j$ on $D^n_{i-j=d}$, substituting into
$l^{n,p}_{i,j} = (i + p(j)) \bmod n$ separates $s$ and $d$ from the column \emph{step functions}
\[
    s^{sum}_{p,j} := (p(j) - j) \bmod n
    \;\;\text{and}\;\;
    s^{dif\!f}_{p,j} := (p(j) + j) \bmod n,
\]
giving $l^{n,p}_{s-j,j} = (s + s^{sum}_{p,j}) \bmod n$ and
$l^{n,p}_{d+j,j} = (d + s^{dif\!f}_{p,j}) \bmod n$. Thus $l^{n,p}$ is injective along $D^n_{i+j=s}$
(resp.\ $D^n_{i-j=d}$) iff $j \mapsto s^{sum}_{p,j}$ (resp.\ $j \mapsto s^{dif\!f}_{p,j}$) is
injective on the columns meeting that diagonal. When $\gcd(n,6) > 1$, step collisions are unavoidable
(Section~\ref{sec:combined}); our majorants $f^+_f$ (sum) and $f^-_f$ (diff) lower one colliding
column into a vacant step residue so Lemma~\ref{lem:inj-maj} applies. Within $F_f$, we write
$\val(i,j) := l^{n,p_f}_{i,j}$, $s^+_j := s^{sum}_{p_f,j}$, and $s^-_j := s^{dif\!f}_{p_f,j}$.

\subsection[Family 1]{Family $F_1$\,: $n = 1,5 \bmod 6$, $n \ge 5$}

Define $p_1(0) := 0$, $p_1(j) := n - 2j$ for $0 < 2j < n$, and $p_1(j) := 2n - 2j$ for $2j \ge n$, so
$p_1(j) = (-2j) \bmod n$. Since $\gcd(2,n) = \gcd(3,n) = 1$, $p_1$, $s^+_j = (-3j) \bmod n$, and
$s^-_j = (-j) \bmod n$ are bijections on $[0,n-1]$. Setting $f^+_1 := f^-_1 := \val$ gives:

\begin{lemma}\label{lem:fam1}
$L^{n,p_1}$ is an incremental cyclic Latin square.
\end{lemma}

\begin{proof}
Immediate from Lemma~\ref{lem:inj-maj} since $f^+_1 = f^-_1 = \val$ and $s^+_j, s^-_j$ are bijections
on $[0,n-1]$.
\end{proof}

\subsection[Family 2]{Family $F_2$\,: $n = 2,4 \bmod 6$, $n \ge 10$}

Here $n = 2m$ ($m \ge 5$, $\gcd(3,n) = 1$). Define $p_2(j) := (2j + m + 1) \bmod n$ on $[0,m-1]$ and
$p_2(j) := (2(j-m) + m + 4) \bmod n$ on $[m,n-1]$. Since $2m = n$, the two blocks step by $+2$
through opposite parities $(m+1) \bmod 2$ and $(m+4) \bmod 2$, so $p_2$ is a permutation of $[0,n-1]$,
with step functions $s^+_j = (j + m + 1) \bmod n$, $s^-_j = (3j + m + 1) \bmod n$ on $[0,m-1]$ and
$s^+_j = (j - m + 4) \bmod n$, $s^-_j = (3j - m + 4) \bmod n$ on $[m,n-1]$ (injective on each block
since $\gcd(3,n) = 1$).

\emph{Collisions and majorants.}
For $j_1 \in [0,m-1]$ and $j_2 \in [m,n-1]$ ($1 \le j_2 - j_1 \le n - 1$, $2m = n$):
\begin{itemize}[nosep]
  \item \emph{Sum collision derivation}: For $j_1 \in [0,m-1]$ and $j_2 \in [m,n-1]$, the sum steps are
  $s^+_{j_1} = (j_1 + m + 1) \bmod n$ and $s^+_{j_2} = (j_2 - m + 4) \bmod n$. A collision occurs iff
  $s^+_{j_1} = s^+_{j_2}$, which means $(j_2 - m + 4) - (j_1 + m + 1) \equiv 0 \pmod n$, or equivalently
  $(j_2 - j_1 - 2m + 3) \equiv 0 \pmod n$. Since $2m = n$, this simplifies to $(j_2 - j_1 + 3) \equiv 0 \pmod n$.
  Because $j_1 \le m - 1$ and $j_2 \ge m$, their difference satisfies $1 \le j_2 - j_1 \le n - 1$, and therefore
  $4 \le j_2 - j_1 + 3 \le n + 2$. Since $n \ge 10$, the only multiple of $n$ in the range $[4, n + 2]$ is $n$
  itself. Hence $j_2 - j_1 + 3 = n$, uniquely giving $j_2 - j_1 = n - 3$.
  Along any sum diagonal $D^n_{i+j=s}$, cells satisfy $i_1 + j_1 = i_2 + j_2 = s$, so
  $i_1 - i_2 = j_2 - j_1 = n - 3$. Since $i_1, i_2 \in [0,n-1]$, we must have $i_1 \ge n - 3$ and $i_2 \le 2$;
  likewise $j_2 \ge n - 3$ and $j_1 \le 2$. Thus, collisions occur exclusively between the $3 \times 3$
  bottom-left corner $(i_1 \in [n-3,n-1], j_1 \in [0,2])$ and the $3 \times 3$ top-right corner
  $(i_2 \in [0,2], j_2 \in [n-3,n-1])$. Shifting the top-right $3 \times 3$ corner defines
  $f^+_2(i,j) := j - (n - 3) \in \{0, 1, 2\}$ for $i < 3, \; j \ge n - 3$, and $\val(i,j)$ otherwise.
  \item \emph{Diff collision derivation}: The diff steps are $s^-_{j_1} = (3j_1 + m + 1) \bmod n$ and
  $s^-_{j_2} = (3j_2 - m + 4) \bmod n$. Then $s^-_{j_2} - s^-_{j_1} \equiv 3(j_2 - j_1) - 2m + 3 \equiv 3(j_2 - j_1 + 1) \pmod n$
  (using $2m = n$). Since $\gcd(3,n) = 1$ when $n \equiv 2,4 \pmod 6$, $3(j_2 - j_1 + 1) \equiv 0 \pmod n \iff
  j_2 - j_1 + 1 \equiv 0 \pmod n$. With $1 \le j_2 - j_1 \le n - 1$, we have $2 \le j_2 - j_1 + 1 \le n$, whose only
  multiple of $n$ is $n$, forcing $j_2 - j_1 = n - 1$. This uniquely selects $j_1 = 0$ and $j_2 = n - 1$.
  Along diff diagonal $i_2 - i_1 = j_2 - j_1 = n - 1$, this forces $(i_1,j_1) = (0,0)$ and $(i_2,j_2) = (n-1,n-1)$
  on $d = 0$. Shifting $(n-1,n-1)$ defines $f^-_2(n-1,n-1) := 1$, and $\val(i,j)$ otherwise.
\end{itemize}

\begin{table}[htbp]
\centering
\footnotesize
\begin{tabular}{@{}lllll@{}}
\toprule
Dir. & Shifted region & Unshifted $\val(i,j)$ & Majorant $f_2(i,j)$ & Step offset on diagonal \\
\midrule
Sum & $i \in [0,2], \; j = n - 3 + l$ ($l \in [0,2]$) & $i + m - 2 + 2l \ge 3$ & $l \in [0,2]$ & $(2l - i - n + 3) \bmod n$ \\
Diff & $(i,j) = (n-1,n-1)$ ($d = 0$) & $m + 1 \ge 6$ & $1$ & $1 \notin \{s^-_j \mid j < n-1\}$ \\
\bottomrule
\end{tabular}
\caption{Collision regions and majorant shifts for Family~$F_2$ ($n = 2m \ge 10$).}
\label{tab:fam2}
\end{table}

\begin{lemma}\label{lem:fam2}
$L^{n,p_2}$ is an incremental cyclic Latin square.
\end{lemma}

\begin{proof}
We verify Lemma~\ref{lem:inj-maj} (Table~\ref{tab:fam2}) using $2m = n$.
\emph{Sum (1)}: For $i \in [0,2]$ and $j = n - 3 + l \ge m$ ($l \in [0,2]$),
$2(n - 3 + l - m) + m + 4 = n + m - 2 + 2l$ ($2m = n$) and $3 \le m - 2 \le i + m - 2 + 2l \le m + 4 < n$,
so $\val(i,n-3+l) = i + m - 2 + 2l \ge 3 > l = f^+_2(i,n-3+l) \ge 0$.
\emph{Sum (2)}: On $D^n_{i+j=s}$, shifted cells have $f^+_2(s-j_2,j_2) = j_2 - (n-3)$ (injective in
$j_2$), and unshifted cells have distinct $\val$-values. For unshifted $(i_1,j_1)$ and shifted
$(i_2,j_2)$ ($i_2 \in [0,2]$, $i_1 = i_2 + j_2 - j_1$), $(\val(i_1,j_1) - f^+_2(i_2,j_2)) \bmod n$
reduces (via $2m = n$) to $i_2 + j_1 + m - 2 \in [m-2, n-1]$ if $j_1 \in [0,m-1]$ and to
$i_2 + (j_1 - m) + 1 \in [1, m+2]$ if $j_1 \in [m,n-1]$; both lie in $[1,n-1]$ (no collision).
\emph{Diff (1)}: At $(n-1,n-1)$, $n - 1 + 2(n - 1 - m) + m + 4 = 2n + m + 1$ ($2m = n$, $6 \le m + 1 < n$),
so $\val(n-1,n-1) = m + 1 > 1 = f^-_2(n-1,n-1) > 0$.
\emph{Diff (2)}: On $d = 0$, unshifted $(j,j)$ ($j \le n-2$) are distinct, and $\val(j,j) \ne 1$
(using $2m = n$ and $\gcd(3,n) = 1$):
$(3j + m + 1) \bmod n = 1 \iff 3(j - m) \bmod n = 0 \iff j = m \notin [0,m-1]$, and
$(3j - m + 4) \bmod n = 1 \iff 3(j + 1 - m) \bmod n = 0 \iff j = m - 1 \notin [m,n-2]$.
\end{proof}

\subsection[Family 3]{Family $F_3$\,: $n = 3 \bmod 6$, $n \ge 27$}

Here $n = 3k$ ($k \ge 9$, $\gcd(2,n) = \gcd(4,n) = 1$). On the blocks $I_u := [uk, (u+1)k-1]$
($u \in \{0,1,2\}$), define $p_3(j) := (3j + 2k + 4u) \bmod n$. Since $3k = n$ and $\gcd(4,3) = 1$,
the blocks $I_u$ step by $+3$ through the residue classes $(2k + 4u) \bmod 3$, so $p_3$ is a
permutation of $[0,n-1]$, with step functions $s^+_j = (2j + 2k + 4u) \bmod n$ and
$s^-_j = (4j + 2k + 4u) \bmod n$ on $I_u$ (injective on each $I_u$ since $\gcd(2,n) = \gcd(4,n) = 1$).

\emph{Collisions and majorants.}
For $j_1 \in I_{u_1}$ and $j_2 \in I_{u_2}$ ($0 \le u_1 < u_2 \le 2$):
\begin{itemize}[nosep]
  \item \emph{Sum collision derivation}: The sum steps satisfy $s^+_{j_2} - s^+_{j_1} \equiv 2(j_2 - j_1 + 2(u_2 - u_1)) \pmod n$.
  Since $\gcd(2,n) = 1$, this vanishes iff $(j_2 - j_1 + 2(u_2 - u_1)) \equiv 0 \pmod n$.
  If $u_2 - u_1 = 1$, then $1 \le j_2 - j_1 \le 2k - 1$, so $3 \le j_2 - j_1 + 2 \le 2k + 1 < 3k = n$, which has no
  multiple of $n$.
  If $u_2 - u_1 = 2$ (so $u_1 = 0, u_2 = 2$), then $k + 1 \le j_2 - j_1 \le 3k - 1 = n - 1$, giving
  $k + 5 \le j_2 - j_1 + 4 \le n + 3$. Since $k \ge 9$, the only multiple of $n$ in this interval is $n$, forcing
  $j_2 - j_1 + 4 = n$, i.e., $j_2 - j_1 = n - 4$.
  Along $i_1 + j_1 = i_2 + j_2$, $i_1 - i_2 = j_2 - j_1 = n - 4$ forces $j_1, i_2 \in [0,3]$ and $i_1, j_2 \in [n-4,n-1]$.
  Shifting the top-right $4 \times 4$ corner defines $f^+_3(i,j) := \val(i,j) - k$ ($i < 4, \; j \in \{n-4,n-3\}$),
  $\val(i,j) - 2k$ ($i < 4, \; j \in \{n-2,n-1\}$), and $\val(i,j)$ otherwise.
  \item \emph{Diff collision derivation}: The diff steps satisfy $s^-_{j_2} - s^-_{j_1} \equiv 4(j_2 - j_1 + u_2 - u_1) \pmod n$.
  Since $\gcd(4,n) = 1$, this vanishes iff $(j_2 - j_1 + u_2 - u_1) \equiv 0 \pmod n$.
  If $u_2 - u_1 = 1$, then $2 \le j_2 - j_1 + 1 \le 2k < n$, with no multiple of $n$.
  If $u_2 - u_1 = 2$ ($u_1 = 0, u_2 = 2$), then $k + 3 \le j_2 - j_1 + 2 \le n + 1$, whose only multiple of $n$ is $n$
  (since $k \ge 9$). This forces $j_2 - j_1 = n - 2$.
  On $i_2 - i_1 = j_2 - j_1 = n - 2$, this forces $i_1, j_1 \in \{0,1\}$ and $i_2, j_2 \in \{n-2,n-1\}$.
  Shifting the bottom-right $2 \times 2$ corner defines $f^-_3(i,n-2) := \val(i,n-2) - 2k$ and
  $f^-_3(i,n-1) := \val(i,n-1) - k$ ($i \ge n - 2$), and $\val(i,j)$ otherwise.
\end{itemize}

\begin{table}[htbp]
\centering
\footnotesize
\begin{tabular}{@{}lllll@{}}
\toprule
Dir. & Shifted corner column & Unshifted $\val(i,j)$ & Majorant $f_3(i,j)$ & Offset $v_l$ \\
\midrule
Sum & $i \in [0,3], \; j = n - 4 + l, \; l \in \{0,1\}$ & $i + 2k - 4 + 3l \ge 2k - 4$ & $\val(i,j) - k$ & $v_0 = k, \; v_1 = k + 2$ \\
Sum & $i \in [0,3], \; j = n - 4 + l, \; l \in \{2,3\}$ & $i + 2k - 4 + 3l \ge 2k + 2$ & $\val(i,j) - 2k$ & $v_2 = 4, \; v_3 = 6$ \\
Diff & $i \in [n-2,n-1], \; j = n - 2$ ($r = i - n + 2$) & $r + 2k \ge 2k$ & $r$ & $0$ \\
Diff & $i \in [n-2,n-1], \; j = n - 1$ ($r = i - n + 2$) & $r + 2k + 3 > k$ & $r + k + 3$ & $k + 4$ \\
\bottomrule
\end{tabular}
\caption{Corner collision regions and majorant shifts for Family~$F_3$ ($n = 3k \ge 27$).}
\label{tab:fam3}
\end{table}

\begin{lemma}\label{lem:fam3}
$L^{n,p_3}$ is an incremental cyclic Latin square.
\end{lemma}

\begin{proof}
We verify Lemma~\ref{lem:inj-maj} (Table~\ref{tab:fam3}) using $3k = n$.
\emph{Sum (1)}: For $i \in [0,3]$ and $j = n - 4 + l \in I_2$ ($l \in [0,3]$), reduction modulo~$n$
($3k = n$) gives $\val(i,n-4+l) = i + 2k - 4 + 3l \in [2k-4, 2k+8]$, exceeding $k$ ($l \in \{0,1\}$)
and $2k$ ($l \in \{2,3\}$), so $0 < f^+_3(i,j) < \val(i,j)$.
\emph{Sum (2)}: On $D^n_{i+j=s}$, $f^+_3(i,n-4+l) = (s + v_l) \bmod n$ with distinct
$v_0 = k, v_1 = k + 2, v_2 = 4, v_3 = 6$ ($k \ge 9$). None equals $s^+_{j_1}$ ($j_1 \in I_{u_1}$):
using $3k = n$ and $\gcd(2,n) = 1$, $s^+_{j_1} = v_l \bmod n$ reduces to $j_1 = k - 2u_1 + l \notin I_{u_1}$
($l \in \{0,1\}$) and $j_1 = 2k - 2u_1 + l \notin I_{u_1}$ ($l \in \{2,3\}$).
\emph{Diff (1)}: For $i = n - 2 + r$ and $j = n - 2 + l$ ($r, l \in \{0,1\}$), reduction modulo~$n$
($3k = n$) gives $\val(i,j) = r + 2k + 3l \in [2k, 2k+4]$, so $f^-_3(i,n-2) = r \ge 0$ and
$f^-_3(i,n-1) = r + k + 3 > 0$.
\emph{Diff (2)}: On $D^n_{i-j=d}$, $r = d + l$ gives shifted values $(d + 0) \bmod n$ and
$(d + k + 4) \bmod n$. Using $3k = n$ and $\gcd(4,n) = 1$, $s^-_{j_1} = 0$ and
$s^-_{j_1} = (k + 4) \bmod n$ reduce to $j_1 = k - u_1 \notin I_{u_1}$ and
$j_1 = 2k + 1 - u_1 \notin I_{u_1}$.
\end{proof}

\subsection[Family 4]{Family $F_4$\,: $n = 0 \bmod 6$, $n \ge 12$}

Here $n = 2m = 3k$ ($m \ge 6, \; k \ge 4$). Define $p_4(0) := n - 2$, $p_4(j) := 2j + 2$ on
$[1, m-3]$, $p_4(m-2) := 0$, $p_4(j) := 2(j - m + 1) + 1$ on $[m-1, n-2]$, and $p_4(n-1) := 2$. The
pins $\{0, m-2, n-1\}$ cover $\{n-2, 0, 2\}$, $[1,m-3]$ covers $\{4, 6, \dots, n-4\}$, and
$[m-1,n-2]$ covers $\{1, 3, \dots, n-1\}$, so $p_4$ is a permutation of $[0,n-1]$. Using $2m = n$,
its step functions $(s^+_j, s^-_j)$ modulo~$n$ are $(n-2, n-2)$ at $j = 0$, $(j+2, 3j+2)$ on
$[1,m-3]$, $(m+2, m-2)$ at $j = m-2$, $(j+3, 3j+3)$ on $[m-1,n-2]$, and $(3, 1)$ at $j = n-1$.

\emph{Collisions and majorants.}
\begin{itemize}[nosep]
  \item \emph{Sum}: $s^+_j$ omits $\{2, m, m+1\}$ and repeats $m + 2$ at $\{m-2, m-1\}$ ($i_2 < n - 1$),
  $n - 2$ at $\{0, n-5\}$ ($i_2 < 5$, i.e.\ $i_2 + j_2 < n$), and $3$ at $\{1, n-1\}$ ($i_2 < 2$).
  Shifting these seams into $\{m+1, m, 2\}$ defines $f^+_4(i,m-1) := \val(i,m-1) - 1$ ($i < n-1$),
  $f^+_4(i,n-5) := \val(i,n-5) - (m-2)$ ($i + j < n$), $f^+_4(i,n-1) := \val(i,n-1) - 1$ ($i < 2$),
  and $\val(i,j)$ otherwise.
  \item \emph{Diff collision derivation}: Since $n = 3k$, the diff steps satisfy $s^-_j \equiv 3j + 2 \equiv 2 \pmod 3$
  on the first block $[1,m-3]$, and $s^-_j \equiv 3j + 3 \equiv 0 \pmod 3$ on the second block $[m-1,n-2]$.
  Hence the residue class $d + 1 \pmod 3$ is completely vacant across both unshifted blocks.
  Within $[1,m-3]$, $s^-_{j_1} = s^-_{j_2} \iff 3(j_2 - j_1) \equiv 0 \pmod n \iff k \mid (j_2 - j_1)$.
  Since the block width is $m - 3 < 2k$, the only multiple is $j_2 - j_1 = k$.
  Along $i_2 - i_1 = j_2 - j_1 = k$, this forces $j_2 \ge 1 + k$ and $i_2 \ge k$, defining
  Region~A ($1+k \le j \le m-3, \; i \ge k$).
  Similarly, within $[m-1,n-2]$, $j_2 - j_1 = k$ forces $j_2 \ge m - 1 + k$ and $i_2 \ge k$.
  The column $j = m - 1 + k$ forms the boundary between blocks: for $k \le i < 2k$ it is Region~B,
  and for $i \ge 2k$ it is Region~C (with $n = 12$ vs.\ $n \ge 18$ split). Columns $m + k \le j \le n - 2$
  form Region~D. Shifting Regions A, B, D by $-1$ or $-2$ moves them into the vacant residue $d + 1 \pmod 3$.
  Specifically, $f^-_4(i,j) := \val(i,j) - 1$ on Region~A and Region~B, $\val(i,j) - 4$ on Region~C$_{12}$
  ($j = 9, \; i \ge 8, \; n = 12$), $\val(i,j) - 5$ on Region~C$_{\ge 18}$ ($j = m-1+k, \; i \ge 2k, \; n \ge 18$),
  $\val(i,j) - 2$ on Region~D ($m+k \le j \le n-2, \; i \ge k$), and $\val(i,j)$ otherwise.
\end{itemize}

\begin{table}[htbp]
\centering
\footnotesize
\begin{tabular}{@{}lllll@{}}
\toprule
Dir. & Shifted region / pin & Unshifted $\val(i,j)$ & Majorant $f_4(i,j)$ & Offset on diagonal \\
\midrule
Sum & Seams $j = m - 1, \; n - 5, \; n - 1$ & $i+1, \; i+n-7, \; i+2$ & $i, \; i+m-5, \; i+1$ & $m+1, \; m, \; 2 \pmod n$ \\
Diff & Region A ($1+k \le j \le m-3, \; i \ge k$) & $i + 2j + 2 - n \ge 4$ & $\val(i,j) - 1$ & $[4, m-8]$ \\
Diff & Region B ($j = m-1+k, \; k \le i < 2k$) & $i - k + 1 \ge 1$ & $\val(i,j) - 1$ & $m - 1$ \\
Diff & Region C ($j = m-1+k, \; i \ge 2k$) & $i - k + 1 \ge k + 1$ & $\val - 4$ ($n=12$), $\val - 5$ & $2$ ($n=12$), $m - 5$ \\
Diff & Region D ($m+k \le j \le n-2, \; i \ge k$) & $i + 2j - 2n + 3 \ge 3$ & $\val(i,j) - 2$ & $[m+1, n-5]$ \\
Diff & Pins $j = 0, \; m - 2, \; n - 1$ & $(i + p_4(j)) \bmod n$ & $\val(i,j)$ & $-2, \; m - 2, \; n + 1$ \\
\bottomrule
\end{tabular}
\caption{Seam and block collision regions and offsets for Family~$F_4$ ($n = 2m = 3k \ge 12$).}
\label{tab:fam4}
\end{table}

\begin{lemma}\label{lem:fam4}
$L^{n,p_4}$ is an incremental cyclic Latin square.
\end{lemma}

\begin{proof}
We verify Lemma~\ref{lem:inj-maj} (Table~\ref{tab:fam4}) using $2m = 3k = n$.
\emph{Sum (1) \& (2)}: On the three seams, reduction modulo~$n$ ($2m = n$) gives $\val(i,m-1) = i + 1 \ge 1$,
$\val(i,n-5) = i + n - 7 > m - 2$, and $\val(i,n-1) = i + 2 \ge 2$, so $0 \le f^+_4(i,j) < \val(i,j)$,
and on $D^n_{i+j=s}$ the shifted seams take the vacant step values $\{m+1, m, 2\}$, making $f^+_4$
injective.
\emph{Diff (1)}: Reduction modulo~$n$ gives $\val(i,j) = i + 2j + 2 - n \ge 4$ on~A,
$\val(i,j) = i - k + 1 \ge 1$ on~B ($\ge 5$ on C$_{12}$, $\ge 7$ on C$_{\ge 18}$), and
$\val(i,j) = i + 2j - 2n + 3 \ge 3$ on~D, so $0 \le f^-_4(i,j) < \val(i,j)$.
\emph{Diff (2)}: On $D^n_{i-j=d}$, the unshifted parts of $[1,m-3]$ and $[m-1,n-2]$ are injective in
$d + 2$ and $d \pmod 3$. Substituting $i = d + j$, Regions A--D and pin $j = m - 2$ lie in
$d + 1 \pmod 3$, while pins $j = 0$ and $j = n - 1$ have offsets $-2$ ($i \ge 2$; for $i \in \{0,1\}$,
$\val(i,0) \in \{n-2,n-1\}$ exceeds all other values) and $n + 1$ ($i \le n - 3$; for
$i \in \{n-2,n-1\}$, $\val(i,n-1) \in \{0,1\}$ is below all other values). Using $3k = 2m$, the
offsets $f^-_4(i,j) - d$ satisfy
\[
    -2 < 4 \le 3(j - k) + 1 \le m - 8 < m - 5 < m - 2 < m - 1 < m + 1 \le 3(j - m - k) + m + 1 \le n - 5 < n + 1
\]
for $n \ge 18$ (and $-2 < 2 < 4 < 5 < 7 < 13$ for $n = 12$, where Region~A is empty), separating all
pieces.
\end{proof}

\subsection[Existence of Incremental Cyclic Latin Squares]{Existence of Incremental Cyclic Latin Squares}

Families $F_1$--$F_4$ cover all $n \ge 4$ ($n \notin \{4, 6\}$) except $n \in \{8, 9, 15, 21\}$ below
$n \ge 10$ ($F_2$) and $n \ge 27$ ($F_3$). Verifying these four small cases is a finite check
(mechanized in Lean~4~\cite{PalomarLean,CompanionRepo}); explicit witnesses are in
Appendix~\ref{app:small} (Lemma~\ref{lem:small}).

\begin{theorem}\label{thm:main2}
For every $n \in \mathbb{N}$ with $n \ge 4$ and $n \notin \{4, 6\}$, there exists an
\emph{incremental} cyclic Latin square $L^{n,p}$.
\end{theorem}

\begin{proof}
Partition $n \ge 4$ ($n \notin \{4, 6\}$) by $n \bmod 6$ and apply Lemma~\ref{lem:fam1}
($n = 1, 5 \bmod 6$), Lemma~\ref{lem:small} or \ref{lem:fam2} ($n = 2, 4 \bmod 6$),
Lemma~\ref{lem:small} or \ref{lem:fam3} ($n = 3 \bmod 6$), and Lemma~\ref{lem:fam4}
($n = 0 \bmod 6, \; n \ge 12$).
\end{proof}

\section[Combined Existence, Obstructions, and Complexity]{Combined Existence, Obstructions, and Complexity}\label{sec:combined}

\emph{Why no cyclic solution exists for $n \in \{4, 6\}$.}
When $\gcd(n,6) > 1$, the Hall--Paige / toroidal $n$-queens obstruction shows that no permutation $p$
of $\mathbb{Z}_n$ makes both $j \mapsto (p(j) \mp j) \bmod n$ bijections (summing $p(j) - j$ modulo~$n$
gives a parity contradiction when $2 \mid n$, and similarly when $3 \mid n$). Thus every $L^{n,p}$ with
$\gcd(n,6) > 1$ has step collisions. For $n \ge 8$, the board is large enough to lower colliding
corners/seams into vacant step residues; for $n \in \{4, 6\}$, exhaustive search over all $4! = 24$ and
$6! = 720$ permutations confirms that every $L^{n,p}$ violates
$|\,D^{n,k}\,| \le k + 1$~\cite{CompanionRepo}. Explicit non-cyclic incremental Latin squares $L^4, L^6$
(\cite{PalomarLean,CompanionRepo}) are in Appendix~\ref{app:noncyclic} (Lemma~\ref{lem:L4L6}).

\emph{Algorithmic complexity.}
Evaluating any cell $l^{n,p}_{i,j} = (i + p(j)) \bmod n$ or querying coin membership
$(i,j) \in S^n_c \iff l^{n,p}_{i,j} < c$ requires:
\begin{enumerate}[nosep]
  \item At most $3$ integer comparisons and $2$ modular additions to evaluate $p(j)$ ($O(1)$ operations);
  \item $1$ integer addition and $1$ modulo-$n$ reduction to compute $l^{n,p}_{i,j}$ ($O(1)$ operations);
  \item $1$ integer comparison to check $l^{n,p}_{i,j} < c$ ($O(1)$ operations).
\end{enumerate}
Thus, querying any entry takes \textbf{$O(1)$ time}.
Generating or materializing the \textbf{entire $n \!\times\! n$ Latin square} takes $n^2 \times O(1) = \mathbf{O(n^2)}$ time,
while storing the closed-form construction requires only \textbf{$O(n)$ space} for the permutation $p$.
This improves upon Le, Li and Wang~\cite{LLW1990} by an \textbf{order of $n$}: because~\cite{LLW1990}
constructs solutions separately for each $c$, generating all $n+1$ solutions requires
$\sum_{c=1}^n O(n \cdot c) = \mathbf{\Theta(n^3)}$ coin placements.
Finally, as shown in Section~\ref{sec:inj-maj}, injective majorization replaces the $\mathbf{O(n^3)}$
operations required by naive diagonal prefix verification ($(4n - 2) \text{ diagonals} \times n \text{ thresholds} \times O(n) \text{ cells}$)
by a single symbolic check.

\begin{theorem}\label{thm:main3}
For every $n \in \mathbb{N}$ with $n \ge 4$, there exists an \emph{incremental} Latin square $L^n$.
\end{theorem}

\begin{proof}
Combine Theorem~\ref{thm:main2} ($n \notin \{4,6\}$) with Lemma~\ref{lem:L4L6} in
Appendix~\ref{app:noncyclic} ($n \in \{4,6\}$)~\cite{PalomarLean,CompanionRepo}.
\end{proof}

\begin{theorem}\label{thm:main4}
For every $n \in \mathbb{N}$ with $n \ge 4$, there exists an \emph{incremental} Latin square $L^n$
representing $n + 1$ solutions of $(n,c)$-coins problems \emph{simultaneously} for all $c \in [0,n]$
as $S^n_c = \{L^n_{i,j} \mid l^n_{i,j} < c\}$.
\end{theorem}

\begin{proof}
Immediate from Theorems~\ref{thm:main1} and \ref{thm:main3}.
\end{proof}

\subsection{Lean 4 Formalization Architecture and Coverage}\label{sec:formalization}

All definitions, lemmas, and theorems in this paper are formalized in Lean~4 with
\textbf{zero \texttt{sorry}s, zero \texttt{native\_decide}s, and zero linter warnings}~\cite{PalomarLean,CompanionRepo}.
Table~\ref{tab:lean-coverage} details the exact correspondence between the paper and the formal codebase.

\begin{table}[htbp]
\centering
\scriptsize
\begin{tabular}{@{}llll@{}}
\toprule
Paper Result & Lean 4 Declaration & Proof Method & Formal File \\
\midrule
Def.~\ref{def:nc} ($(n,c)$-coins) & \texttt{IsNCCoinsSolution} & Formal Definition & \texttt{IncrementalNCCoins.lean} \\
Def.~\ref{def:is} (Incremental) & \texttt{IsIncremental} & Formal Definition & \texttt{IncrementalNCCoins.lean} \\
Def.~\ref{def:cs} (Cyclic square) & \texttt{cyclicSquare} & Formal Definition & \texttt{IncrementalNCCoins.lean} \\
Thm.~\ref{thm:main1} (Simultaneity) & \texttt{incremental\_represents\_nc\_coins} & Constructive proof & \texttt{IncrementalNCCoins.lean} \\
Lem.~\ref{lem:inj-maj} (Inj.\ Maj.) & \texttt{val\_le\_of\_inj\_le} & Constructive proof & \texttt{IncrementalNCCoins.lean} \\
Lem.~\ref{lem:fam1} ($F_1$) & \texttt{solve\_family1} & Constructive (\texttt{omega}) & \texttt{Solution.lean} \\
Lem.~\ref{lem:fam2} ($F_2$) & \texttt{solve\_family2} & Constructive (\texttt{omega}) & \texttt{Solution.lean} \\
Lem.~\ref{lem:fam3} ($F_3$) & \texttt{solve\_family3} & Constructive (\texttt{omega}) & \texttt{Solution.lean} \\
Lem.~\ref{lem:fam4} ($F_4$) & \texttt{solve\_family4} & Constructive (\texttt{omega}) & \texttt{Solution.lean} \\
Lem.~\ref{lem:small} ($n \in \{8,9,15,21\}$) & \texttt{solve\_8, solve\_9, 15, 21} & Kernel \texttt{decide} & \texttt{Solution.lean} \\
Thm.~\ref{thm:main2} (Cyclic existence) & \texttt{exists\_cyclic\_incremental} & Constructive dispatch & \texttt{Solution.lean} \\
Lem.~\ref{lem:L4L6} ($n \in \{4,6\}$) & \texttt{isIncremental\_L4, L6} & Kernel \texttt{decide} & \texttt{Solution.lean} \\
Thm.~\ref{thm:main3}--\ref{thm:main4} (Main) & \texttt{exists\_incremental\_representing\_nc\_coins} & Constructive proof & \texttt{Solution.lean} \\
\bottomrule
\end{tabular}
\caption{Mapping between paper results and declarations in the Lean~4 formalization.}
\label{tab:lean-coverage}
\end{table}

The infinite families $F_1$--$F_4$ are proved constructively for arbitrary $n$ using a two-tier
\emph{Flat Disjunctive Specification Pattern}: each piecewise permutation $p_f$ and majorant $f^\pm_f$
is proved equivalent to a flat disjunction of linear equations, enabling Lean's \texttt{omega} tactic
to verify diagonal injectivity in under $25$ seconds without exponential case explosion.
For the small dimensions $n \in \{4, 6, 8, 9, 15, 21\}$, diagonal injectivity is verified by Lean's
trusted kernel evaluator (\texttt{decide}) in milliseconds.

\appendix

\section[Small Cyclic Cases]{Small Cyclic Cases $n \in \{8, 9, 15, 21\}$}\label{app:small}

We record the explicit permutations $p_n$ and majorants $f^+, f^-$ (equal to
$\val(i,j) := l^{n,p_n}_{i,j}$ outside the shifted cells listed below) for $n \in \{8, 9, 15, 21\}$
(\cite{PalomarLean,CompanionRepo}):
\begin{itemize}[nosep]\raggedright
  \item \emph{$n = 8$}: $p_8 = (0, 4, 1, 3, 6, 2, 7, 5)$ ($s^+_j$ repeats $0$ at $\{0,3\}$, omitting
  $4$; $s^-_j$ repeats $5$ at $\{1,6\}$, omitting $1$). Shifted cells: $f^+(i,3) = i + 1$ ($i < 4$),
  $f^+(4,3) = 3$; $f^-(i,6) = i - 5$ ($i \in \{5,6\}$), $f^-(7,6) = 5$.
  \item \emph{$n = 9$}: $p_9 = (0, 6, 3, 1, 8, 5, 2, 4, 7)$ ($s^-_j$ is a permutation; $s^+_j$ repeats
  at $\{0,5\}, \{1,6\}$, omitting $2, 3$). Shifted cells: $f^+(i,5) = i + 4$ ($i < 3$), $f^+(3,5) = 2$;
  $f^+(i,6) = i$ ($i < 2$), $2i - 3$ ($i \in \{2,3\}$).
  \item \emph{$n = 15$}: $p_{15} = (0, 3, 6, 9, 12, 8, 5, 2, 13, 1, 4, 7, 10, 14, 11)$ ($s^+_j$ is a
  permutation; $s^-_j$ repeats at $\{3,13\}, \{9,14\}$, omitting $2, 5$). Shifted cells:
  $f^-(10,13) = 5$, $f^-(i,13) = i - 8$ ($i \in [11,14]$); $f^-(i,14) = i - 5$ ($i \in \{5,6\}$),
  $i - 6$ ($i \in [7,10]$), $2$ ($i \in \{11,12,14\}$), $1$ ($i = 13$).
  \item \emph{$n = 21$}: $p_{21} = (0, 18, 15, 12, 9, 6, 3, 7, 4, 1, 14, 17, 10, 16, 2, 5, 11, 8, 20, 13, 19)$
  ($s^+_j$ repeats at $\{0,7\}, \{1,8\}, \{2,9\}, \{3,14\}$, omitting $7,8,10,14$; $s^-_j$ repeats at
  $\{2,18\}, \{5,19\}$, omitting $2,5$). Shifted cells: $f^+(i,7) = i + 6$ ($i \le 12$), $13$
  ($i = 13$); $f^+(i,8) = i + 3$ ($i \le 7$), $i + 2$ ($i \in [8,10]$), $i + 1$ ($i \in \{11,13\}$),
  $9$ ($i = 12$); $f^+(i,9) = i$ ($i \le 7$), $i - 2$ ($i \in \{8,9,10,12,13\}$), $7$ ($i = 11$);
  $f^+(i,14) = i + 1$ ($i \in [0,5] \cup [7,9]$), $6$ ($i = 6$); $f^-(i,18) = i - 13$ ($i \in [16,20]$);
  $f^-(i,19) = i - 14$ ($i \in [14,16]$), $i - 17$ ($i \in [17,20]$).
\end{itemize}

\begin{lemma}\label{lem:small}
For each $n \in \{8, 9, 15, 21\}$, $L^{n,p_n}$ is an incremental cyclic Latin square.
\end{lemma}

\begin{proof}
Finite verification of $0 \le f^\pm_n(i,j) \le \val(i,j)$ and diagonal injectivity into the omitted
step residues proves Lemma~\ref{lem:inj-maj} in each case~\cite{PalomarLean,CompanionRepo}.
\end{proof}

\begin{example}[Illustration for $n = 8$]\label{ex:n8}
For $n = 8$, the column permutation is
$p_8 = (0, 4, 1, 3, 6, 2, 7, 5)$.
The step vectors across columns $j \in [0,7]$ are:
\begin{align*}
  s^+ &= (p_8(j) - j) \bmod 8 = (0, 3, 7, \mathbf{0}, 2, 5, 1, 6), \\
  s^- &= (p_8(j) + j) \bmod 8 = (0, \mathbf{5}, 3, 6, 2, 7, \mathbf{5}, 4).
\end{align*}
Here $s^+$ repeats residue $0$ at $\{0,3\}$ while residue $4$ is vacant; $s^-$ repeats residue $5$ at $\{1,6\}$ while residue $1$ is vacant.
Along sum diagonals meeting column $j = 3$, setting $f^+(i,3) := i + 1$ for $i < 4$ and $f^+(4,3) := 3$
shifts the colliding values into the vacant residue $4$ without collisions, satisfying $0 \le f^+ \le \val$.
Similarly, along diff diagonals meeting column $j = 6$, setting $f^-(i,6) := i - 5$ ($i \in \{5,6\}$) and
$f^-(7,6) := 5$ shifts the colliding values into the vacant residue $1$.
Thus, Lemma~\ref{lem:inj-maj} holds constructively.
\end{example}

\section[Non-Cyclic Cases]{Non-Cyclic Cases $n \in \{4, 6\}$}\label{app:noncyclic}

Explicit non-cyclic incremental Latin squares $L^4$ and $L^6$~\cite{PalomarLean,CompanionRepo} are:
\[
    L^4 :=
    \begin{pmatrix}
        2 & 0 & 1 & 3 \\
        1 & 3 & 2 & 0 \\
        0 & 2 & 3 & 1 \\
        3 & 1 & 0 & 2
    \end{pmatrix},
    \qquad
    L^6 :=
    \begin{pmatrix}
        1 & 0 & 2 & 3 & 5 & 4 \\
        3 & 1 & 4 & 0 & 2 & 5 \\
        4 & 3 & 5 & 2 & 1 & 0 \\
        0 & 2 & 1 & 5 & 4 & 3 \\
        2 & 5 & 0 & 4 & 3 & 1 \\
        5 & 4 & 3 & 1 & 0 & 2
    \end{pmatrix}.
\]

\begin{lemma}\label{lem:L4L6}
$L^4$ and $L^6$ are incremental Latin squares.
\end{lemma}

\begin{proof}
A sorted diagonal $(v_0 \le \dots \le v_{r-1})$ satisfies $|\,D^{n,k}\,| \le k + 1$ for all $k$ iff
$v_m \ge m$ ($m < r$). Inspecting the diagonals of $L^4$ (main diagonals sort to $[2,2,3,3]$) and
$L^6$ (main diagonals sort to $[1,2,2,4,5,5]$ and $[1,1,2,3,5,5]$) verifies $v_m \ge m$ on every
diagonal~\cite{PalomarLean,CompanionRepo}.
\end{proof}

\begin{thebibliography}{99}
\raggedright
\bibitem{BS2009}
J. Bell and B. Stevens,
\emph{A survey of known results and research areas for $n$-queens},
Discrete Math.\ \textbf{309} (2009), no.~1, 1--31,
\url{https://www.sciencedirect.com/science/article/pii/S0012365X07010394}.

\bibitem{LLW1990}
M.-H. Le, W. Li, and E. T. Wang,
\emph{A generalization of the $n$-queen problem},
J. Systems Sci.\ Math.\ Sci.\ \textbf{3} (1990), no.~2, 183--192,
\url{https://zbmath.org/4212061}.

\bibitem{Pratt2016}
K. Pratt,
\emph{Closed-form expressions for the $n$-queens problem and related problems},
Int.\ Math.\ Res.\ Not.\ (2019), no.~22, 6978--6988,
\url{https://arxiv.org/abs/1609.09585}.

\bibitem{CW2010}
N. J. Cavenagh and I. M. Wanless,
\emph{On the number of transversals in Cayley tables of cyclic groups},
Discrete Appl.\ Math.\ \textbf{158} (2010), no.~2, 136--146,
\url{https://doi.org/10.1016/j.dam.2009.09.006}.

\bibitem{MY2026}
M. Mirzavaziri and D. Yaqubi,
\emph{Cyclic Latin Eulerian numbers},
arXiv:2609.28808 (2026),
\url{https://arxiv.org/abs/2609.28808}.

\bibitem{CompanionRepo}
M. Duponcheel and L. Duponcheel,
\emph{Companion repository: Incremental $(n,c)$-coins problem},
\url{https://github.com/NofanTasi/incremental-nc-coins-problem}.

\bibitem{PalomarLean}
L. Duponcheel and M. Duponcheel,
\emph{Palomar Lean~4 formalization (\texttt{IncrementalNCCoins.lean}, \texttt{Solution.lean})},
\url{https://github.com/LucDuponcheelAtGitHub/incremental-nc-coins}.

\bibitem{MO510135}
M. Duponcheel, \emph{MathOverflow \#510135}, \url{https://mathoverflow.net/questions/510135}.

\bibitem{PSE138939}
M. Duponcheel, \emph{Puzzling SE \#138939}, \url{https://puzzling.stackexchange.com/questions/138939}.

\bibitem{MO516225}
M. Duponcheel, \emph{MathOverflow \#516225}, \url{https://mathoverflow.net/questions/516225}.
\end{thebibliography}

\end{document}
